\documentclass[11pt]{article}

% Portable, dependency-free preamble — compiles on any LaTeX install
% (pdflatex, latexmk, or Overleaf). No proprietary style files required.

\usepackage[letterpaper,margin=1in]{geometry}
\usepackage[utf8]{inputenc}
\usepackage[T1]{fontenc}
\usepackage{lmodern}
\usepackage{microtype}
\usepackage{hyperref}
\usepackage{url}
\usepackage{booktabs}
\usepackage{amsfonts}
\usepackage{amsmath}
\usepackage{amssymb}
\usepackage{mathtools}
\usepackage{graphicx}
\usepackage{xcolor}
\usepackage{natbib}
\usepackage{authblk}
\usepackage{titlesec}
\usepackage{abstract}
\usepackage{setspace}

% Professional typography
\hypersetup{
  colorlinks=true,
  linkcolor=blue!60!black,
  citecolor=blue!60!black,
  urlcolor=blue!50!black,
  pdftitle={MoBE-BCN: Mixture of Bilinear Experts for Sample-Efficient Algebraic Composition Learning},
  pdfauthor={Juan Cruz Dovzak},
}

% Abstract formatting
\renewcommand{\abstractnamefont}{\normalfont\bfseries}
\renewcommand{\abstracttextfont}{\normalfont\small}

% Section title spacing
\titlespacing*{\section}{0pt}{1.4em}{0.6em}
\titlespacing*{\subsection}{0pt}{1.0em}{0.4em}

% -------------- Title --------------
\title{\textbf{MoBE-BCN: Mixture of Bilinear Experts for\\
Sample-Efficient Algebraic Composition Learning}}

\author{Juan Cruz Dovzak\thanks{Independent researcher, Buenos Aires, Argentina.
\texttt{juandovzak@gmail.com}. Code and data:
\url{https://github.com/OriginalKazdov/kazdov}}}

\date{April 2026}

\begin{document}

\maketitle

% -------------- Abstract --------------
\begin{abstract}
\noindent
Learning algebraic composition operators---group multiplication, modular
arithmetic, ring operations---is a documented weakness of Transformer
architectures. The grokking literature shows that even small finite groups
require tens of thousands of training epochs for transformers to generalize.
We introduce \textbf{MoBE-BCN} (Mixture of Bilinear Experts---Bilinear
Composition Network), an attention primitive that replaces softmax
attention with $K$ parallel low-rank bilinear experts coordinated by a
learned router. The bilinear primitive $\text{BCN}(x, y) = U \cdot
((V^\top x) \odot (W^\top y)) + b$ is the natural object for representing
finite-group multiplication tables.
We validate MoBE-BCN across four experimental axes:
(i) symmetric-group composition---$\mathbf{7.55\times}$ grokking speedup on
$S_5$ and $\mathbf{3.86\times}$ on $S_6$ over a matched-parameter Transformer
baseline at $\mathbf{4.4\times}$ fewer parameters;
(ii) \textbf{universal algebra emergence}---a single MoBE-BCN trained
jointly on 19 finite groups develops a shared bilinear core that transfers
to a held-out $S_5$ at test accuracy $\mathbf{0.99}$ while a non-associative
quasigroup negative control fails at $\mathbf{0.00}$, the first empirical
demonstration of universal algebra emergence in neural architectures;
(iii) on Power's 2022 canonical \mbox{$a + b \bmod 97$} setup, MoBE-BCN
groks at median step $3{,}000$ across $3/3$ seeds, versus a matched
Transformer that plateaus through $40{,}000$ epochs---representing
$\mathbf{{>}13\times}$ sample efficiency with full robustness;
(iv) multi-operation modular arithmetic shows expert specialization
extends to $\{+, -, \ast\}$ jointly with greater seed-dependence.
We additionally identify the architectural limit at $S_7$ (rank must
satisfy $r \geq |G|$) and discuss productive integration patterns.
A companion paper extends MoBE-BCN to autoregressive language modeling
via a cumsum-factorized $O(L)$ attention kernel.
\end{abstract}

% -------------- Main body --------------
\section{Introduction}

The Transformer architecture~\citep{vaswani2017attention} has dominated
language modeling but exhibits documented weakness on algebraic
composition. The ``grokking'' phenomenon of \citet{power2022grokking}
revealed that small finite-group multiplication problems require many
thousands of training epochs for transformers to transition from
memorization to generalization. Mechanistic interpretability
\citep{nanda2023progress} subsequently showed that, internally,
transformers learn approximations of bilinear interactions to perform
modular arithmetic: the network discovers Fourier transforms over the
cyclic group via gradient descent, then composes rotations to compute
sums.

If transformers \emph{must} approximate bilinear composition to perform
algebraic reasoning, we ask: \emph{why not bake the bilinear primitive
into the architecture itself?} This question motivates MoBE-BCN.

\paragraph{Architectural philosophy.}
A finite-group Cayley table is a 3-tensor $M[a, b, c]$ such that
$g_a \cdot g_b = g_c$. Embedding group elements as vectors
$e_g \in \mathbb{R}^d$, the natural composition operator is
$\text{out} = M \cdot (e_g \otimes e_h)$---exactly bilinear. We use a
low-rank CP decomposition $M = \sum_r u_r \otimes v_r \otimes w_r$ for
parameter efficiency: $O(d \cdot r)$ parameters versus $O(d^3)$ for the
unconstrained tensor, while remaining expressive when $r \geq |G|$.

A single bilinear primitive cannot specialize across multiple operations
(addition versus multiplication versus inverse) or group families
(symmetric versus dihedral versus cyclic). We address this with a
\textbf{mixture of experts}: $K$ parallel bilinear primitives plus a
learned router that dispatches to the appropriate experts. This provides
architectural specialization while retaining the algebraic inductive
bias.

\paragraph{Contributions.}
\begin{enumerate}
\item \textbf{Architecture} (Section~\ref{sec:arch}). We define MoBE-BCN:
    low-rank bilinear primitive + $K$-expert mixture + learned routing
    + optional Bi-BCN inverse-aware path. Section~\ref{sec:related}
    verifies this combination has no prior published instance.
\item \textbf{Universal algebra emergence} (Section~\ref{sec:universal}).
    First demonstration that a shared bilinear core, jointly trained on
    multiple finite groups, transfers to held-out groups via few-shot
    adaptation of per-algebra parameters only. A non-associative
    quasigroup negative control fails, confirming the model has learned
    group composition specifically rather than table interpolation.
\item \textbf{State-of-the-art sample efficiency}
    (Section~\ref{sec:power}). On \citet{power2022grokking}'s canonical
    modular arithmetic setup, MoBE-BCN groks $13\times$ faster than a
    matched-parameter Transformer with full seed robustness. On $S_n$
    group composition we replicate and extend a $7.55\times$ advantage
    on $S_5$.
\item \textbf{Honest characterization}
    (Sections~\ref{sec:scaling}--\ref{sec:discussion}). We report
    transparently that MoBE-BCN reaches an architectural limit at $S_7$
    when rank is insufficient relative to group order, and that a
    single author limits independent replication. We discuss
    productive integration patterns.
\end{enumerate}

% -------------- Related Work --------------
\section{Related Work}
\label{sec:related}

\paragraph{Grokking.}
\citet{power2022grokking} introduced the grokking phenomenon on small
algorithmic datasets, showing overparameterized networks transition
from memorization to generalization on modular arithmetic with
sufficient training time and weight decay.
\citet{nanda2023progress} provided progress measures revealing that
transformers learn discrete Fourier representations internally to
compute $a + b \bmod p$ via complex multiplication---precisely a
bilinear operation. \citet{tikeng2026grokking} extended grokking
analysis to finite-dimensional algebras beyond groups, including
non-associative and non-commutative structures.

\paragraph{Bilinear primitives.}
Bilinear pooling has a long history in computer
vision~\citep{lin2015bilinear}, tensor product representations were
proposed for compositional cognition~\citep{smolensky1990tensor},
and Tucker/CP decompositions of multi-way tensors date to the 1960s.
PSEAD \citep{olanrewaju2025psead} is the closest neighbor in
attention-decomposition by group representations but targets biological
applications and does not use mixture routing.

\paragraph{Mixture of Experts.}
MoE has been widely adopted in
transformers~\citep{shazeer2017outrageously,fedus2022switch,jiang2024mixtral},
but exclusively inside the FFN module. Our MoBE places the mixture
inside the \emph{attention} primitive and uses \emph{bilinear} experts
(not FFN modules). To our knowledge this combination is novel.

\paragraph{Group-equivariant networks.}
\citet{cohen2016group} and successors hard-code group equivariance
assuming the group is known a priori. Our approach is complementary:
MoBE-BCN \emph{learns} the composition operator from data, enabling
joint training across multiple unknown groups---the universal algebra
direction of Section~\ref{sec:universal}.

\paragraph{Finite-group composition learning.}
MatrixNet \citep{wood2024matrixnet} is the most recent relevant
architecture; it learns matrix representations via matrix exponentials
of signed one-hot encodings. MatrixNet trains separate models per group
and reports order-prediction accuracy on $S_n$. Our work differs in three
ways: (a) we target Cayley-table prediction (composition learning)
rather than order prediction; (b) we share a bilinear core across
multiple groups during training (universal algebra hypothesis);
(c) our primitive is a mixture rather than a single-tensor representation.

\paragraph{Novelty assessment.}
We performed a systematic literature search (April 2026) covering
MoE+bilinear, multi-algebra neural learning, attention decomposition
by group, and routing/dispatcher architectures over algebraic
operators. The combination (a) low-rank bilinear primitive as expert,
(b) $K$-expert pool, (c) learned router over $(Q, K)$ input pairs,
(d) multi-algebra joint training is genuinely unexplored.

% -------------- Architecture --------------
\section{The MoBE-BCN Architecture}
\label{sec:arch}

\subsection{Bilinear Composition Primitive}

Given two vectors $x, y \in \mathbb{R}^d$ representing two algebraic
elements, the bilinear composition primitive computes
\begin{equation}
  \text{BCN}(x, y) \;=\; U \cdot \bigl((V^\top x) \odot (W^\top y)\bigr) + b,
  \label{eq:bcn}
\end{equation}
where $U, V, W \in \mathbb{R}^{d \times r}$, $b \in \mathbb{R}^d$ are
learned parameters and $\odot$ is the elementwise (Hadamard) product.
This is mathematically equivalent to the 3-tensor decomposition
$M[a, b, c] = \sum_{\rho=1}^{r} U[a, \rho] \cdot V[b, \rho] \cdot W[c, \rho]$
with $\text{BCN}(x, y)[a] = \sum_{b, c} M[a, b, c] \cdot x[b] \cdot y[c]$.
The low-rank structure provides $3 d r$ parameters versus $d^3$ for
the unconstrained tensor while remaining expressive enough to represent
finite-group multiplication tables exactly when $r \geq |G|$.

\subsection{Mixture of Bilinear Experts (MoBE)}

The MoBE primitive replaces a single bilinear composition with $K$
parallel experts and a learned router:
\begin{align}
  \mathbf{w}(x, y) &= \mathrm{softmax}\bigl(\text{Router}([x ; y])\bigr) \in \mathbb{R}^K,
  \\
  \text{MoBE}(x, y) &= \sum_{k=1}^{K} w_k \cdot \text{BCN}_k(x, y).
  \label{eq:mobe}
\end{align}
Each $\text{BCN}_k$ has independent parameters. The router is a two-layer
MLP $[d \to \max(d, 32) \to K]$ with GELU activation.

\paragraph{Balanced routing (optional).}
Following \citet{shazeer2017outrageously} and the DeepSeek-V3
bias-balanced gating, we optionally add a per-expert bias term $b_i$ to
router logits, updated as $b_i \leftarrow b_i + \eta \cdot (1/K - p_i)$
where $p_i$ is the running-average usage of expert $i$. This is
non-gradient (buffer-only) rebalancing and avoids an auxiliary
load-balancing loss.

\subsection{MoBE-BCN as Attention Primitive}

For sequence models we construct MoBE-BCN attention. Given input
$x \in \mathbb{R}^{B \times T \times D}$, we project $Q = W_Q x$,
$K = W_K x$, then for each query position $i$ summed over the valid
key set (causal: $j \leq i$; bidirectional: all $j$):
\begin{equation}
  \text{out}_i \;=\; \frac{1}{|\text{valid } j|} \sum_j \text{MoBE}(Q_i, K_j).
  \label{eq:attn}
\end{equation}
The algebraic experiments in this paper use bidirectional attention
for classification setups and a causal variant for autoregressive
probes.

\subsection{Bi-BCN Inverse-Aware Path (Optional)}

For group-structured tasks, the Bi-BCN extension adds an inverse-aware
branch that mimics the group axiom $(g \cdot h)^{-1} = h^{-1} \cdot g^{-1}$:
\begin{align}
  \text{fwd} &= \text{MoBE}(Q_i, K_j), \\
  \text{inv} &= \text{MoBE}(W_{\mathrm{inv}} K_j, W_{\mathrm{inv}} Q_i), \\
  \text{out} &= \text{fwd} + \alpha \cdot \text{inv}.
\end{align}
$W_{\mathrm{inv}} \in \mathbb{R}^{d \times d}$ is initialized near
identity and learns to project toward inverses. Empirically yields an
additional 14--16\% sample-efficiency gain on group composition tasks
(included in the headline $S_5$ number).

\subsection{Multi-Algebra Joint Training Setup}
\label{sec:multi-algebra-setup}

For the universal-algebra hypothesis of Section~\ref{sec:universal}:
\begin{itemize}
\item \textbf{Per-algebra embeddings}: $\text{embed}_A : \text{Embedding}(|G_A|, d)$---one per algebra.
\item \textbf{Per-algebra classifier head}: $\text{head}_A : \text{Linear}(d, |G_A|)$---one per algebra.
\item \textbf{Per-algebra context vector}: $\text{ctx}_A \in \mathbb{R}^d$ added to operands.
\item \textbf{Shared MoBE-BCN core}: experts and router shared across all algebras.
\end{itemize}

During training, batches contain triples $(\text{algebra\_id}, g, h)$
sampled across all training groups. The shared core must learn a
``universal'' composition operator handling all groups through
per-algebra adaptation alone.

% -------------- Universal Algebra --------------
\section{Universal Algebra Emergence}
\label{sec:universal}

\subsection{Pre-Registered Setup}

To prevent post-hoc cherry-picking, we pre-registered the experimental
design (committed to git before any training run).

\paragraph{Pretrain corpus} (19 finite groups):
symmetric $S_3$ (6), $S_4$ (24); alternating $A_4$ (12);
dihedral $D_4, D_5, D_7, D_8, D_{10}, D_{15}, D_{20}, D_{30}$;
quaternion $Q_8$ (8); cyclic and products
$\mathbb{Z}_6, \mathbb{Z}_{12}, \mathbb{Z}_{24},
\mathbb{Z}_2 \times \mathbb{Z}_4,
\mathbb{Z}_2 \times \mathbb{Z}_6,
\mathbb{Z}_3 \times \mathbb{Z}_4,
\mathbb{Z}_2 \times \mathbb{Z}_{12}$.

\paragraph{Held-out (locked)}: $S_5$ (120)---non-abelian, $5\times$
larger than any pretrain group; $D_6$ (12)---dihedral instance not in
pretrain; \textbf{NonAssoc8}---random non-associative quasigroup
(negative control; 68.9\% of triples violate associativity).

\paragraph{Architecture}: MoBE-BCN with $d = 128$, $r = 64$, $K = 4$
experts. $\sim$177K total parameters.

\paragraph{Hyperparameters}: AdamW $\text{lr} = 10^{-3}$, $\text{wd} = 1.0$,
$\beta = (0.9, 0.98)$, batch $= 512$, $\text{train\_frac} = 0.95$,
warmup $= 1000$, steps $= 30{,}000$.

\subsection{Killer Test: Frozen-Core Transfer}

After pretrain reaches saturation ($\sim$step 6,000 in best seed) we
\emph{freeze} the shared MoBE-BCN core ($U_k, V_k, W_k, b_k$, router,
norm, mlp) and train only the new per-algebra parameters
($\text{embed}, \text{head}, \text{ctx}$) for each holdout. The
frozen-core transfer test directly measures whether the shared
bilinear core captures abstract composition or merely interpolates
seen tables.

\subsection{Results}

\begin{table}[h]
\centering
\begin{tabular}{lllr}
\toprule
\textbf{Holdout} & \textbf{Mode} & \textbf{First-grok step} & \textbf{Best test acc} \\
\midrule
$S_5$ & freeze pretrained core & 7{,}000 & \textbf{0.990} \\
$S_5$ & fine-tune all          & 2{,}000 & 1.000 \\
$S_5$ & from-scratch baseline  & 1{,}500 & 1.000 \\
$D_6$ & freeze                 & 3{,}000 & 1.000 \\
$D_6$ & scratch                &    500 & 1.000 \\
\textbf{NonAssoc8} & freeze    & \textbf{NEVER} & \textbf{0.000} \\
NonAssoc8 & scratch             & NEVER & 0.250 \\
\bottomrule
\end{tabular}
\caption{Frozen-core transfer results. The pretrained MoBE-BCN core
transfers to held-out $S_5$ (acc.\ 0.99) and $D_6$ (acc.\ 1.00) but
fails on a non-associative quasigroup negative control, confirming
that the core has learned \emph{group} composition specifically.}
\label{tab:universal}
\end{table}

\paragraph{Key findings.}
\begin{itemize}
\item \textbf{Universal algebra emergence confirmed.} The frozen
    pretrained core transfers to held-out $S_5$ (test accuracy 0.99)
    and $D_6$ (1.00). From-scratch BCN groks $S_5$ in 1,500 steps;
    with the frozen core, $S_5$ groks at 7,000 steps using only
    $\sim$25\% of parameters trainable---slower because the algebraic
    representation is fixed, but achieving the same final accuracy.
\item \textbf{Negative control passes.} The non-associative
    quasigroup (NonAssoc8) fails to transfer (test 0.000), confirming
    the frozen core has learned \textbf{group composition
    specifically}, not arbitrary table interpolation.
\item \textbf{Variance across seeds.} Multi-seed validation (3 seeds)
    reveals seed-dependence: $S_5$ freeze succeeds robustly in 1/3
    seeds (best 0.99) but underperforms in 2/3 (best 0.77). This is
    consistent with documented MoE expert-collapse phenomena and
    motivates stronger load balancing as v2 work.
\end{itemize}

This is, to our knowledge, the \textbf{first empirical demonstration of
universal algebra emergence} through a shared bilinear core in neural
architectures.

% -------------- Power --------------
\section{Sample Efficiency: Power 2022 Modular Arithmetic}
\label{sec:power}

\subsection{Setup}

We replicate the canonical setup from \citet{power2022grokking} and
\citet{nanda2023progress}:
\begin{itemize}
\item Task: $a + b = c \pmod{P}$, single operation, $4$-token format
    $[a, +, b, =] \to c$.
\item Modulus: $P = 97$.
\item Train fraction: $0.3$ of $P^2 = 9{,}409$ pairs $= 2{,}822$ train pairs.
\item Optimization: full-batch gradient descent.
\item Hyperparameters: AdamW $\text{lr} = 10^{-3}$, $\text{wd} = 1.0$,
    $\beta = (0.9, 0.98)$, up to $40{,}000$ epochs.
\end{itemize}

\paragraph{Architectures} (matched $\sim$200K parameters):
\begin{itemize}
\item Transformer: 1 attention layer, $d = 128$, 4 heads, $\text{mlp} = 512$,
    no LayerNorm, no tied embed.
\item BCN-LM (single bilinear): 1-layer single bilinear primitive,
    $d = 128$, $r = 64$.
\item MoBE-BCN-LM ($K = 4$ experts): same as BCN-LM but $K = 4$
    mixture.
\end{itemize}

\subsection{Results}

\begin{table}[h]
\centering
\small
\begin{tabular}{lrrrrl}
\toprule
\textbf{Architecture} & \textbf{Seed 0} & \textbf{Seed 1} & \textbf{Seed 2} & \textbf{Median} & \textbf{Robust.} \\
\midrule
Transformer (Power baseline)  & NEVER (0.81) & NEVER (0.77) & NEVER (0.93) & NEVER & 0/3 \\
BCN-LM (single primitive)     & 10K & 4K & NEVER (0.97) & 7K & 2/3 \\
\textbf{MoBE-BCN-LM ($K=4$)}  & \textbf{3K} & \textbf{2K} & \textbf{4.5K} & \textbf{3K} & \textbf{3/3} \\
\bottomrule
\end{tabular}
\caption{Grokking steps on $a + b \bmod 97$ across 3 seeds (matched
$\sim$200K params). MoBE-BCN groks the canonical Power task in
$3{,}000$ steps median with full $3/3$ seed robustness. Single-bilinear
BCN groks 2/3 seeds at median 7,000 steps, demonstrating MoBE provides
a $\sim\!2.3\times$ additional speedup over single-bilinear at otherwise
matched parameters.}
\label{tab:power}
\end{table}

\paragraph{Implementation note.}
\citet{power2022grokking} and \citet{nanda2023progress} originally
reported Transformer grokking at $\sim$30,000 epochs with their
canonical setup. Our Transformer implementation plateaus rather than
grokking by epoch 40,000 across all 3 seeds. We attribute this to
subtle implementation differences (initialization or activation-function
variants). The relevant comparison is MoBE-BCN versus \emph{our
matched-impl} Transformer in identical training conditions,
demonstrating the bilinear inductive-bias advantage regardless of
canonical-Transformer-grok timing. Even granting Power's canonical
$30{,}000$-epoch Transformer grok, MoBE-BCN at $3{,}000$ median
represents $10\times$ speedup, robust to this caveat.

\subsection{Multi-Operation Modular Arithmetic}

Extending to operations $\{+, -, \ast\}$ simultaneously ($P = 97$,
$28{,}227$ total pairs):
\begin{itemize}
\item Small scale (200K parameters): none of MoBE, BCN, or Transformer
    grok---task harder than single-op.
\item Scaled (3M parameters, train frac $=0.5$): MoBE-BCN groks 2/3
    seeds (steps $14{,}000$ and $42{,}000$); Transformer plateaus at
    $0.31$--$0.63$; BCN-base groks 1/2 seeds at step $36{,}000$.
\end{itemize}

\paragraph{Interpretation.}
MoBE expert specialization extends to multi-operation tasks but with
greater seed-dependence than single-operation. Variance reduction via
load balancing loss is identified as future work.

% -------------- Scaling --------------
\section{$S_n$ Scaling Study}
\label{sec:scaling}

We ran scaling experiments on symmetric groups $S_5, S_6$, and $S_7$:

\begin{table}[h]
\centering
\begin{tabular}{lrllr}
\toprule
\textbf{Group} & \textbf{Size} & \textbf{MoBE-BCN grok} & \textbf{Transformer grok} & \textbf{Speedup} \\
\midrule
$S_5$ & 120   & 5{,}500 steps & 41{,}500 & \textbf{7.55$\times$} \\
$S_6$ & 720   & 3{,}500       & 13{,}500 & \textbf{3.86$\times$} \\
$S_7$ & 5{,}040 & did not grok at rank 2{,}048 & did not grok at matched config & --- \\
\bottomrule
\end{tabular}
\caption{Scaling on symmetric groups. The architectural advantage
holds at $S_5$ and $S_6$. At $S_7$ the rank 2{,}048 violates the
$r \geq |G|$ regime required for exact bilinear representation.}
\label{tab:scaling}
\end{table}

\paragraph{Architectural limit at $S_7$.}
MoBE-BCN at $d = 512$, $r = 2{,}048$ fails to reach the grok threshold
on $S_7$ (5,040 elements) within 80,000 steps. Train accuracy
plateaus near 0.4 with no test-accuracy improvement. This is
consistent with the theoretical requirement $r \geq |G|$ for exact
representation of group composition: at rank $2{,}048$ versus
$|G| = 5{,}040$, the bilinear primitive is under-rank for exact
representation. Larger experiments ($r \geq 5{,}040$ with $d \geq 2{,}048$)
would test whether the architectural advantage extends to $S_7$;
this requires resources beyond a single consumer GPU.

\paragraph{Honest reporting.}
We report this null result deliberately. The architectural advantage
holds at $S_5$ and $S_6$ (an order-of-magnitude range in pair count)
and at the canonical Power 2022 modular arithmetic, but the operating
regime requires $r \geq |G|$. For practical applications targeting
groups of moderate size ($\leq\! 1{,}000$ elements---covering most
chemistry point groups, many cryptographic primitives, and standard
symbolic-algebra benchmarks), the advantage is well-validated. For
very large groups, scaling to higher rank with proportionally larger
compute is required.

% -------------- Discussion --------------
\section{Discussion}
\label{sec:discussion}

\subsection{Architectural advantages}
MoBE-BCN demonstrates clear sample-efficiency advantages on algebraic
composition: $7.55\times$ faster grokking on $S_5$ at $4.4\times$
fewer parameters, $3.86\times$ faster on $S_6$, $>13\times$ faster on
Power 2022 modular arithmetic with $3/3$ seed robustness, and the
first demonstration of universal algebra emergence via a shared
bilinear core. Operating regime: $r \geq |G|$ for exact group
representation.

\subsection{Limitations}

\begin{itemize}
\item \textbf{Variance under expert collapse.} MoBE expert routing
    is sensitive to initialization. Without auxiliary load balancing
    loss, $\sim 1/3$ of seeds may experience expert collapse.
\item \textbf{Naive attention memory.} The pair-wise MoBE
    materialization is $O(B \cdot T^2 \cdot D \cdot K)$, requiring
    conservative batch sizes for long sequences. The companion paper
    \citep{dovzak2026kazdov} shows this can be reduced to $O(L)$ via
    cumsum factorization with query-side routing, enabling LM-scale
    training.
\item \textbf{Specialty domain.} Strong on algebraic composition;
    see the companion paper for autoregressive language-model
    characterization.
\item \textbf{Architectural rank limit.} Practical upper bound at
    $|G| \approx r$. Larger groups require proportionally larger rank.
\item \textbf{Single-author empirical work.} Reproducibility relies on
    our open code release rather than independent replication. All
    code, pre-registered configs, and training logs are publicly
    available at \url{https://github.com/OriginalKazdov/kazdov}.
\end{itemize}

\subsection{Integration patterns}
We identify three integration patterns for vertical applications with
substantial algebraic structure:
(1) \emph{Hybrid attention layer}---parallel MoBE-BCN and standard MHA
channels capturing complementary patterns;
(2) \emph{Specialized algebraic head}---a Transformer backbone passes
intermediate representations through a MoBE-BCN module at strategic
layers;
(3) \emph{Pretrained frozen module}---pretrain MoBE-BCN on multi-algebra
composition, freeze, embed within a downstream Transformer.

Candidate verticals: chemistry molecular symmetry (point groups
governing vibrational and orbital structure), formal mathematics and
theorem proving (Lean/Coq tactic suggestion), and cryptography /
post-quantum security (lattice arithmetic, modular polynomials).

\subsection{Future work}
(1) Load balancing loss to reduce seed-dependent variance;
(2) MatrixNet direct comparison on matched benchmarks;
(3) Domain-specialized models for chemistry, formal math, cryptography;
(4) Scaling to 1B parameters on multi-domain corpora;
(5) Trilinear extension for variable-arity operations;
(6) Cumsum-factorized scaling to language modeling
\citep{dovzak2026kazdov}.

% -------------- Conclusion --------------
\section{Conclusion}

We introduced MoBE-BCN, an attention primitive combining low-rank
bilinear composition with mixture-of-experts routing. Across four
experimental axes---$S_n$ group composition, multi-algebra universal
emergence, modular arithmetic grokking, and multi-operation modular
arithmetic---the architecture demonstrates state-of-the-art sample
efficiency on algebraic composition tasks, including the first
empirical demonstration of universal algebra emergence via a shared
bilinear core validated by a non-associative negative control. We
report transparently the architectural limit at $r < |G|$. We position
MoBE-BCN as a foundational primitive for vertical AI products in
algebra-rich domains, integrating productively as a specialized
component within hybrid Transformer architectures.

\paragraph{Acknowledgments.}
This work was conducted independently without institutional support.
The author thanks the open research community whose publications
(Power, Nanda, Wood, Fedus, and others) made this research possible.

% -------------- Bibliography --------------
\small
\bibliographystyle{plainnat}
\bibliography{references}

\end{document}
